\section{Cyclic decomposition and determinantal representations}
The cyclic projectors
\[
P_3=\tfrac14(I+S^3+S^6+S^9),\qquad
P_4=\tfrac13(I+S^4+S^8)
\]
have images of dimensions three and four. If $E=P_3-P_4$, then
\begin{equation}
\frac{\Tr E}{12}=-\frac1{12}.
\label{eq:minus12}
\end{equation}
The scalar comes from a normalized rank difference, not from an
ordinary sum of positive integers.

\begin{theorem}[Determinantal representation of the channel readout]
Let $\mathcal H_3=\operatorname{Fix}(S^3)$ and
$\mathcal H_4=\operatorname{Fix}(S^4)$. For $0<s<1$,
\begin{equation}
F(s)=\frac{\det_{\mathcal H_3}(I-sS)}{\det_{\mathcal H_4}(I-sS)}
=\frac{1-s^3}{1-s^4}
=\frac{1+s+s^2}{1+s+s^2+s^3}.
\label{eq:det}
\end{equation}
Its continuous extension at $s=1$ equals $3/4$.
\end{theorem}
\begin{proof}
The restriction of $S$ to each image is a cycle of three or four
positions, respectively. Their characteristic polynomials give
$\det(I-sS)=1-s^r$ for $r=3,4$. Cancelling the common factor $1-s$
leaves the final expression, which is regular at $s=1$.
\end{proof}

The substitution $s=q^{30}$ expresses
$F(q^{30})=(1-q^{90})/(1-q^{120})$.
The same pair of spaces that produces \eqref{eq:minus12}
produces the channel quotient. This equality does not identify
a twelve-coordinate matrix with the entire TPK state.

\subsection{Oriented information and the complex sector}
In the antipodal sector $V_-=\operatorname{Ran}((I-S^6)/2)$,
the operator $J=S^3$ satisfies $J^2=-I$ and $J^{\mathsf T}=-J$.
With $Q=P_4(I-S^6)/2$, we have
\[
\operatorname{rank}Q=2,\qquad E|_{V_-}=-Q,\qquad B|_{V_-}=5I-3Q.
\]
On $\operatorname{Ran}Q$, $S^2=-I$ and the determinant is $1+s^2$.
Thus,
\[
F(s)=\frac{1+s+s^2}{1+s}\,\frac1{1+s^2}.
\]
Reversing orientation changes the sign of the form
$\omega(u,v)=\langle Ju,v\rangle$, but preserves $B$ and $F$.
Scalar evaluation does not distinguish all oriented transports.

The transported trace also produces
\[
a_n^{\rm tr}=\Tr(S^nE)=3\mathbf1_{3\mid n}-4\mathbf1_{4\mid n},
\]
and, for $|s|<1$,
\[
-\log F(s)=\sum_{n\ge1}\frac{a_n^{\rm tr}}n\,s^n.
\]
This collection of moments connects the present block with the collected
spectral readouts. Its complete arithmetic development, including
Catalan, will have a separate location; it is not assumed here that an
identity of functions by itself determines its arithmetic nature.
